\section*{PPSZ and width-bounded proof heuristics (Introduction)}
The PPSZ algorithm for $k$-SAT by Paturi, Pudl\'ak, Saks, and Zane [7] is simple to state but famously difficult to analyze. Given a $k$-CNF formula $\Phi$ as input, it first chooses a random ordering $\pi$ of its variables $x_1,\ldots,x_m$. It goes through them one by one, in the order given by $\pi$. For each variable $x$, it tries to derive the correct value using a certain proof heuristic $P$. $P$ takes as input a $k$-CNF formula $\Phi$ and a variable $x$ and returns a value in $\{0,1,?\}$. $P$ must be sound, meaning if $P(\Phi,x)=b\in\{0,1\}$ then $\Phi\models(x=b)$, i.e., every satisfying assignment of $\Phi$ sets $x$ to $b$; however, we allow $P$ to be incomplete, i.e., it may answer "?", meaning "I don't know". If $P(\Phi,x)=b\in\{0,1\}$, then PPSZ sets $x$ to $b$; otherwise it sets $x$ to some $b\in\{0,1\}$ chosen uniformly at random. In either case, it simplifies $\Phi$ to $\Phi|_{x\mapsto b}$. Once all variables have been processed, the resulting formula either contains the empty clause $\square$, and we declare this run of PPSZ a failure; or it does not, in which case PPSZ has found a satisfying assignment.
If PPSZ has success probability $p$ then we can repeat it $1/p$ times, obtaining a constant success probability. As long as $P$ runs in subexponential time, the overall running time of this Monte Carlo algorithm is dominated by $1/p$ (which will, most likely, be exponential in $n$). Which proof heuristics $P$ should one consider? There are currently just two on the market. The first one is $P_w$, which checks whether $(x=b)$ is implied by a set of up to $w$ clauses of $\Phi$. The second one is $R_w$, which tries to derive the clause $(x=b)$ by resolution, bounded by width $w$. Obviously they both can be implemented in time $O^*(\binom{|\Phi|}{w})\le O^*(\binom{m^k}{w})$, which is subexponential as long as $w\in o(m/\log m)$. It is easy to see that $R_{w\cdot k}$ is at least as strong as $P_w$. We also speak of weak PPSZ when it uses $P_w$ and strong PPSZ when it uses $R_w$ (ignoring the concrete values of $w$).
\section*{2.1 Notation}
Given a set of variables $X$, a partial assignment is a function $\alpha:X\to\{0,1,*\}$, that is an assignment of 0-1 values to some of the variables with $*$ intended to mean unset by $\alpha$. We denote the set of variables to which $\alpha$ gives a value by $\operatorname{var}(\alpha):=\{x\in X:\alpha(x)\in\{0,1\}\}$. For two partial assignments $\alpha$ and $\beta$ we write $\alpha\subseteq\beta$ to mean that for every $x\in\operatorname{var}(\alpha)$, it holds that $\beta(x)=\alpha(x)$. Naturally, $\alpha\subset\beta$ means that $\alpha\subseteq\beta$ and $|\operatorname{var}(\alpha)|<|\operatorname{var}(\beta)|$. Given a variable $x$ and $b\in\{0,1\}$, $x\mapsto b$ is the partial assignment which sets $x$ to $b$. The assignment which sets every variable to0 is denoted by $\mathbf0$. For $Y\subseteq X$, we write $Y\mapsto0$ to denote the partial assignment which sets all variables in $Y$ to0. If $\operatorname{var}(\alpha)\cap\operatorname{var}(\beta)=\varnothing$, we define $\alpha\cup\beta$ to be the partial assignment which sets all $x\in\operatorname{var}(\alpha)$ to $\alpha(x)$, all $x\in\operatorname{var}(\beta)$ to $\beta(x)$, and all other variables to $*$. Finally the restriction of a formula $\Phi$ by $\alpha$ is denoted by $\Phi|_\alpha$.
